\documentclass{article}

% Symbols
%\usepackage{recycle}
\usepackage{amsfonts, amsthm}
\usepackage{upgreek}
\usepackage{physics}
\usepackage{cancel}
\usepackage{amssymb, latexsym, amsmath}

% Proof
\renewcommand*{\proofname}{\textbf{Demostraci\'on:}}

% Graphics
\usepackage{graphicx}
\usepackage{pgf}

% Color a letras.
%\usepackage[usenames,dvipsnames,svgnames,table]{xcolor}

% Tikz
\usepackage{tikz}
\usetikzlibrary{arrows,automata}
\usepackage{tikz}
\usetikzlibrary{arrows,automata}

\usetikzlibrary{shapes,calc}
\tikzstyle{edge}=[shorten <=2pt, shorten >=2pt,
  >=stealth, line width=1.1pt]
\tikzstyle{blueE}=[shorten <=2pt, shorten >=2pt,
  >=stealth, line width=1.5pt, blue]
\tikzstyle{blackV}=[circle, fill=black,
  minimum size=6pt,
  inner sep=0pt, outer sep=0pt]
\tikzstyle{blueV}=[circle, fill=blue, draw,
  minimum size=6pt, line width=0.75pt,
  inner sep=0pt, outer sep=0pt]
\tikzstyle{redV}=[circle, fill=red, draw,
  minimum size=6pt, line width=0.75pt,
  inner sep=0pt, outer sep=0pt]
\tikzstyle{redSV}=[semicircle, fill=red, minimum
  size=3pt, inner sep=0pt, outer sep=0pt,
  rotate=225]
\tikzstyle{blueSV}=[semicircle, fill=blue, minimum
  size=3pt, inner sep=0pt, outer sep=0pt,
  rotate=225]
\tikzstyle{blackSV}=[semicircle, fill=black, minimum
  size=3pt, inner sep=0pt, outer sep=0pt,
  rotate=225]
\tikzstyle{vertex}=[circle, draw, minimum size=6pt,
  line width=0.75pt, inner sep=0pt,
  outer sep=0pt]

% Margins
\addtolength{\voffset}{-1.5cm}
\addtolength{\hoffset}{-1.5cm}
\addtolength{\textwidth}{3cm}
\addtolength{\textheight}{3cm}

%Header-Footer
\usepackage{fancyhdr}
\renewcommand{\headrulewidth}{1pt}

\newcommand{\set}[1]{%
  \left\{ #1 \right\}%
}

%\pagenumbering{gobble} -- Este comando
%                       -- quita el número de página.
\footskip = 50pt
\renewcommand{\headrulewidth}{1pt}

\pagestyle{fancyplain}

\begin{document}
\title{UNIVERSIDAD AUT\'ONOMA DE M\'EXICO\\ Facultad de Ciencias}
\author{Autores:
  \\ Fernanda Villaf\'an Flores
  \\ Fernando Alvarado Palacios
  \\ Adri\'an Aguilera Moreno}
\date{}
\maketitle
\begin{center}
  \includegraphics[scale=0.20]{../Imagen/Portada.jpg}\\[0.4cm]
  \Large
  \bf{Gr\'aficas y Juegos}
  \normalsize
\end{center}
\newpage
\fancyhead[r]{ Gr\'aficas y Juegos 2022-1}
%%%%%%%%%%%%%%%%%%%%%%%%%%%%%%%%%%%%%%%%%%%%%%%%%%%%%
\section*{\LARGE{Tarea 4}}

\begin{enumerate}
  %%%%%%%%%%%%%%%%%%%%%%%%%%%%%%%%%%%% Ejercicio 1. %%%%%%%%%%%%%%%%%%%%%%%%%%%%%
  \item Sea $G$ una gr\'afica no trivial.   Demuestre
    que $G$ es una trayectoria si y s\'olo si $G$ es
    un \'arbol con exactamente dos v\'ertices de
    grado $1$.
    
    %%%%%%%%%%%%%%%%%%%%%%%%%%%%%%%%%%%% Ejercicio 2. %%%%%%%%%%%%%%%%%%%%%%%%%%%
  \item \begin{enumerate}
    %~~~~~~~~~~~~~~~~~~~~~~~~~~~~~~~~~~~~~ 2 (a).
    \item Demuestre que cada \'arbol con grado m\'aximo
      $\Delta > 1$ tiene al menos $\Delta$ hojas.
      %~~~~~~~~~~~~~~~~~~~~~~~~~~~~~~~~~~~ 2 (c).
    \item Construya, para cada elecci\'on de $n$ y $\Delta$,
      con $2\le \Delta < n$, un \'arbol de orden $n$ con
      exactamente $\Delta$ hojas.
  \end{enumerate}
    
    %%%%%%%%%%%%%%%%%%%%%%%%%%%%%%%%%%%% Ejercicio 3. %%%%%%%%%%%%%%%%%%%%%%%%%%%%
  \item Un {\em centro} en una gr\'afica es un v\'ertice
    $u$ tal que $\max_{v \in V} d(u, v)$ es m\'inima.
    Demuestre que un \'arbol tiene exactamente un centro
    o dos centros adyacentes.
    
    %%%%%%%%%%%%%%%%%%%%%%%%%%%%%%%%%%%% Ejercicio 4. %%%%%%%%%%%%%%%%%%%%%%%%%%%%
  \item Demuestre o brinde un contraejemplo: Toda
    gr\'afica con menos aristas que v\'ertices tiene
    una componente que es un \'arbol.
    
    %%%%%%%%%%%%%%%%%%%%%%%%%%%%%%%%%%%% Ejercicio 5. %%%%%%%%%%%%%%%%%%%%%%%%%%%%
  \item Un {\em hidrocarburo saturado} es una
    mol\'ecula $C_mH_n$ en la que cada \'atomo de
    carbono tiene cuatro enlaces, cada
    \'atomo de hidr\'ogeno tiene un enlace, y
    ninguna sucesi\'on de enlaces forma un ciclo.
    Demuestre que para cualquier entero positivo
    $m$, la mol\'ecula $C_mH_n$ existe s\'olo si
    $n = 2m + 2$.
    
    %%%%%%%%%%%%%%%%%%%%%%%%%%%%%%%%%%%% Ejercicio 6. %%%%%%%%%%%%%%%%%%%%%%%%%%%%
  \item Demuestre que una sucesi\'on $(d_1, \dots,
    d_n)$ de enteros positivos es la sucesi\'on de
    grados de un \'arbol si y s\'olo si
    $\sum_{i=1}^n d_i = 2(n-1)$.
\end{enumerate}
  %%%%%%%%%%%%%%%%%%%%%%%%%%%%%%%%%%%% Extras %%%%%%%%%%%%%%%%%%%%%%%%%%%%%%%%
\section*{Puntos Extra}
\begin{enumerate}
  %-------------------------------------- Extra 1.
  \item Para una gr\'afica conexa $G$ definimos la
    gr\'afica de \'arboles de $G$, $\mathcal{T}_G$,
    como la gr\'afica que tiene por v\'ertices a
    todos los \'arboles generadores de $G$, y tal que,
    si $S, T \in V_{\mathcal{T}_G}$, entonces $ST$ es
    una arista de $\mathcal{T}_G$ si y s\'olo si
    existen aristas $e \in E_S - E_T$ y $f \in E_T -
    E_S$ tales que $(S - e) + f = T$.   Demuestre que
    $\mathcal{T}_G$ es conexa.
    %-------------------------------------- Extra 2.
  \item Sea $T$ un \'arbol arbitrario con $k+1$ v\'ertices.
    Demuestre que si $G$ es simple y $\delta \ge k$,
    entonces $G$ tiene una subgr\'afica isomorfa a $T$.
    %-------------------------------------- Extra 3.
  \item Sea $\mathcal{T}$ una familia de sub\'arboles
    de un \'arbol $T$.   Deduzca, por inducci\'on sobre
    $|\mathcal{T}|$, que si cualesquiera dos elementos
    de $\mathcal{T}$ tienen un v\'ertice en com\'un,
    entonces hay un v\'ertice de $T$ que est\'a en
    todos los elementos de $\mathcal{T}$.
      %-------------------------------------- Extra 4.
  \item \begin{enumerate}
    \item Determine todos los arboles $T$ tales que
      $\overline T$ tambi\'en es un \'arbol.

    \item Determine todas las gr\'aficas de orden al
      menos cuatro tales que la subgr\'afica inducida por
      cualesquiera tres de sus v\'ertices es un \'arbol.
    \end{enumerate}
    Justifique detalladamente sus respuestas.

\end{enumerate}

\end{document}
